\subsubsection{Interpolation Grid}
\label{SEC:Files:Interpolation_Grid}

%%%%%%%%%%%%%%%%%%%%%%%%%%%%%%%%%%%%%%%%%%%%%%%%%%%%%%%%%%%%%%%%%%%%%%%%%%%%%%%%

It is sometimes desirable to interpolate the solution from one
grid file onto another grid file.
%
For example, after making modifications to an existing grid or
when performing a formal grid sensitivity study.
%
Initializing the starting solution in this way can alleviate some
of the startup transients that typically occur and can reduce the
overall simulation cost.
%
To perform this task, read the existing grid and solution
(or restart) files as one normally would.  
%
In addition, set the various inputs in the
\inptext{nml\_interpolation} namelist block (\sref{SEC:nml_interpolation})
to read the recipient grid and establish the proper number
of solution points on that grid.
%
These inputs are analogous to and support the full range of
grid formats used during normal grid import.

%%%%%%%%%%%%%%%%%%%%%%%%%%%%%%%%%%%%%%%%%%%%%%%%%%%%%%%%%%%%%%%%%%%%%%%%%%%%%%%%

%% Developer note on how parallel interpolation works.

% Parallel Interpolation
% \begin{itemize}
% \setlength\topsep{0.0em}
% \setlength\itemsep{0.0em}
% \setlength\parsep{0.0em}
% \item
% Read the recipient grid.  Allocate the recipient solution.
% \item
% Move the recipient grid and solution from the normal (\verb|*|)
% data structures to the corresponding grid2 (\verb|g2_*|) locations,
% so that the normal data structures can be re-used for the donor
% information.
% \item
% Read the donor grid.  Read donor solution/restart.
% \end{itemize}
% %
% %The donor grid and solution will be read just as with a normal run.
% %Once the recipient grid is imported, 

%%%%%%%%%%%%%%%%%%%%%%%%%%%%%%%%%%%%%%%%%%%%%%%%%%%%%%%%%%%%%%%%%%%%%%%%%%%%%%%%

After the interpolation is complete, the global filter will be applied to the
entire solution if it is enabled in the input file.
%
This can help \quoterm{soften} the result of the interpolation by removing
unwanted oscillations that might have been added to the solution.
%
The final step of the interpolation utility is to write restart and output
files, after which the program will stop without performing
any iterations.
%
The user should examine the output files using flow visualization
software to look for obvious defects which may have been introduced
into the new solution file.
%
If the interpolated solution is satisfactory, then modify the namelist input file to read
the new grid and restart files, and not perform interpolation again.
%
%% Note:  With version 2+ the standard and interpto\_ inputs are in
%%        different namelist blocks!
%This can easily be accomplished by removing the \inptext{interpto\_}
%prefix of the inputs used during interpolation.

%%%%%%%%%%%%%%%%%%%%%%%%%%%%%%%%%%%%%%%%%%%%%%%%%%%%%%%%%%%%%%%%%%%%%%%%%%%%%%%%

The interpolation utility is one of the few parts of the GFR code that has
OpenMP parallelization.
%
Enabling OpenMP in the \texttt{system\_variables.mk} file can really help improve
the performance of the interpolation algorithm, especially for larger
simulations.

%%%%%%%%%%%%%%%%%%%%%%%%%%%%%%%%%%%%%%%%%%%%%%%%%%%%%%%%%%%%%%%%%%%%%%%%%%%%%%%%

\HIDE{
\vspace{1.0em}%
\begin{itemize}
\setlength\topsep{0.0em}
\setlength\itemsep{0.0em}
\setlength\parsep{0.0em}
\item \raggedright
Is there a recommendation for number of cores to use for parallel
interpolation? (Just a few, same as normal solver run, etc.)
\end{itemize}
} %end HIDE

%%%%%%%%%%%%%%%%%%%%%%%%%%%%%%%%%%%%%%%%%%%%%%%%%%%%%%%%%%%%%%%%%%%%%%%%%%%%%%%%

\SeeAlso{
\inp[nmlInterpolation]{nml\_interpolation} (\sref{SEC:nml_interpolation})
%\inptext{interpto\_gridfile},
%\inptext{interpto\_grid\_format},
%\inptext{interpto\_grid\_scaling\_factor}
%\inptext{interpto\_plot3d\_i\_stride},
%\inptext{interpto\_plot3d\_j\_stride},
%\inptext{interpto\_plot3d\_k\_stride},
%\inptext{interpto\_plot3d\_all\_strides},
%\inptext{interpto\_plot3d\_agglomerate\_order},
%\inptext{interpto\_plot3d\_cutfile},
%\inptext{interpto\_plot3d\_bcsfile},
%\inptext{interpto\_solution\_order}
}

%%%%%%%%%%%%%%%%%%%%%%%%%%%%%%%%%%%%%%%%%%%%%%%%%%%%%%%%%%%%%%%%%%%%%%%%%%%%%%%%

